% !TEX TS-program = lualatexmk
%% %%%%%%%%%%%%%%%%%%%%%%%%%%%%%%%%%%%%%%%%%%%%%%%%%%%%%%%%%%%%%%%
%% Das Masterfile für die Buchfassung: AGFA-Master-Book.tex
%% Stand:	2026-10-01
%% --
%% Unterschiede zu AGFA-Master.tex (scrartcl):
%%   - Klasse scrbook, zweiseitig, Kapitel beginnen rechts
%%   - Gliederung: \chapter (I, II, ...) > \section (1, 2, ...) > \subsection
%%   - Theoreme und Gleichungen: im Text  Theorem 2.1  bzw. (2.1),
%%     im Verweis (\cref, \ref, \eqref)  Theorem I.2.1  bzw. (I.2.1)
%%   - \frontmatter / \mainmatter / \backmatter statt Handarbeit
%%     mit \pagenumbering
%% Die Pakete in ./preamble sind dieselben wie bei AGFA-Master.tex;
%% sie erkennen die Buchklasse selbst.
%% %%%%%%%%%%%%%%%%%%%%%%%%%%%%%%%%%%%%%%%%%%%%%%%%%%%%%%%%%%%%%%%
\documentclass[%	 -- siehe KOMA-Script 
	,toc		= bib 		 
	,parskip	= half-		 
	,headings	= normal	 		 	 		 
%	,headings	= chapterprefix	% "Kapitel I" als eigene Zeile über dem Titel
	,numbers	= noenddot		
	,BCOR 		= 12mm
	,twoside 	= true			% Buch: zweiseitig
	,open		= right			% Kapitel beginnen auf einer rechten Seite
	,cleardoublepage = empty	% eingeschobene Leerseiten ganz leer
	,version 	= last
 	,titlepage	= true   
	,ngerman	
%	,english	
% 	,review					% für das Korrekturlesen % wegmachen	 
	]{scrbook}			 

%% -- Nutze die AGFA-Pakete
%% --
\usepackage[%
%	, libertinus		% font libertinus
%	, lmodern			% font lmodern, das Original von Knuth
%	, thmframed			% wer es gerahmt haben will
	]{./preamble/agfa-art}
	
%% -- Literatur-DB durch eigene ersetzen 
%% --
\addbibresource{./bib/agfa-bib.bib}		% agfa-bib-neu.bib -> abcd-bib.bib	

\ExecuteBibliographyOptions{%
	,backref		= true	% Auf was ist verwiesen; Rest muss weg
  	,url			= true	% false: URL wird nicht angezeigt 
  	,doi			= false	% true: URLDOI wird separat angezeigt
	,eprint			= false	% true: wird separat angezeigt
	,maxcitenames  	= 2 	% viele Autoren -> nur die ersten beiden erscheinen  
    ,maxbibnames   	= 10 	% im Verzeichnis tauchen alle auf
	}

%% -- Links farbig
%% --
\hypersetup{
	,colorlinks	= true    			%   onlineversion true
	,urlcolor	= blue       		%
	,citecolor	= blue      		%
	,linkcolor	= blue			 	% oder black
%  	,hidelinks 						% Vor dem Druck % entfernen
	}
	

%% == Ab hier eigene Angaben ==
	
%% -- Was soll aktuell bearbeitet werden
%% -- Mittels % entsprechend steuern
%% --
\includeonly{%
	./content-book/AGFA-Book-Einleitung,
	./content-book/AGFA-Book-Kapitel-1,
	./content-book/AGFA-Book-Kapitel-2,
	}


%% -- Ausfüllen aber für die finale Version mit % auskommentieren
%% -- (zweiseitig: i = innen, c = Mitte, o = außen)
%% --
\KOMAoptions{footsepline}
\ifoot{Name}
\cfoot{Stand der Arbeit:}
\ofoot{\today}

%% -- Paket comment (\begin{comment} ... \end{comment}) wird bereits
%% -- von agfa-art geladen; siehe Literaturverzeichnis

%% -- Start
%% --
\begin{document}

%% %%%%%%%%%%%%%%%%%%%%%%%%%
%% Frontmatter: römische Seitenzahlen, Kapitel ohne Nummer
%% %%%%%%%%%%%%%%%%%%%%%%%%
\frontmatter

%% -- Bei einer Arbeit zum Schluss bei \begin{comment} und \end{comment}
%% -- ein % hinzufügen und nicht vergessen, titlepage = true setzen
%% --
%\begin{comment}
%%
\begingroup 
	% !TEX root = ../AGFA-Master.tex
%% %%%%%%%%%%%%%%%%%%%%%%%%%%%%%%%
%% ./content/0-AGFA-title.tex
%% Titelseite (frei gestaltet, daher mit titlepage-Umgebung)
%% Stand: 2026-09-30
%% %%%%%%%%%%%%%%%%%%%%%%%%%%%%%%%

\begin{titlepage}
%%
\parbox{.75\linewidth}{%
Universität Tübingen \\ 
Mathematisch-Naturwissenschaftliche Fakultät \\ 
Fachbereich Mathematik} 
\par\vspace{3cm}
%%
\begin{center}\Large\bfseries
	\begin{tabular}{c}
		Bachelor-/Masterarbeit \\[1cm]
		Der Titel der Arbeit \\
		-- Untertitel --
	\end{tabular}%
	\par\normalsize\normalfont\vspace*{2cm}%
%%
	\begin{tabular}{c}
		\large\textbf{Vorname Name} \\[20pt]
		Datum der Abgabe		% beim Schreiben ggf. \today, zur Abgabe festes Datum
	\end{tabular}\par\vspace*{\fill}
%%	
	\begin{tabular}[t]{l}
	Betreuer: \\[1ex]
	Prof.\,Dr. Vorname Name \\
	Prof.\,Dr. Vorname Name		
	\end{tabular}
\end{center}
%%
\end{titlepage}
\cleardoublepage			% Titelseite
	\pagestyle{empty}
	% !TEX root = ../AGFA-Master.tex

%% %%%%%%%%%%%%%%%%%%%%%%%%%%%%%%%
%% ./0-AGFA-EigenErkl.tex
%% Eigenständigkeitserklärung
%% Stand: 2026-09-24
%% %%%%%%%%%%%%%%%%%%%%%%%%%%%%%%%
\thispagestyle{empty}
\section*{Erklärung}
%%
{\RaggedRight
Hiermit erkläre ich, dass ich die vorliegende Arbeit 
\enquote{Titel}
selbstständig verfasst und keine anderen als die angegebenen Quellen und Hilfsmittel benutzt habe.
Alle sinngemäß und wörtlich übernommenen Textstellen aus fremden Quellen wurden kenntlich gemacht.\par}% \par vor } noetig, sonst wirkt \RaggedRight nicht
\vspace{5em}

Tübingen, Datum der Abgabe eintragen		\\[2em]
Vorname Name + Unterschrift

%% Anmerkung: Diese Seite kann auch nach dem Literaturverzeichnis eingereiht werden.
%%
%% Unterschrift nicht vergessen
%%

\cleardoubleoddpage

		% Eigenständigkeitserklärung
	% !TEX root = ../AGFA-Master.tex

%% %%%%%%%%%%%%%%%%%%%%%%%%%%%%%%%
%% ./0-AGFA-thanks.tex
%% Danksagung
%% Stand: 2026-09-24
%% %%%%%%%%%%%%%%%%%%%%%%%%%%%%%%%
\thispagestyle{empty} 
\section*{Danksagung}
Hier kann die Danksagung stehen und warum es diese Arbeit gibt.
Aber: \textbf{maximal eine Seite}.
\cleardoubleoddpage
			% Danksagung
	% !TEX root = ../AGFA-Master.tex

%% %%%%%%%%%%%%%%%%%%%%%%%%%%%%%%%
%% ./0-AGFA-summary.tex
%% Zusammenfassung
%% Stand: 2026-09-24
%% %%%%%%%%%%%%%%%%%%%%%%%%%%%%%%%
\thispagestyle{empty} 
\addsec{Zusammenfassung}
Hier bitte eine kurze Zusammenfassung der Arbeit schreiben.
Diese soll nur die wesentlichen Punkte enthalten.
Auf Definitionen, Sätze \usw kann verzichtet werden, da sie später folgen.
\cleardoubleoddpage
		% Zusammenfassung
	\tableofcontents 
	\thispagestyle{empty}
\endgroup
%%
%\end{comment}

%% %%%%%%%%%%%%%%%%%%%%%%%%%
%% Mainmatter: arabische Seitenzahlen ab 1, Kapitel I, II, ...
%% %%%%%%%%%%%%%%%%%%%%%%%%
\mainmatter

%% -- Jedes Kapitel ist eine eigene Datei mit \chapter{...} am Anfang;
%% -- \include beginnt ohnehin auf einer neuen (rechten) Seite.
%% -- Die Einleitung ist ein Kapitel ohne Nummer: \addchap{Einleitung}
%% --
%% %%%%%%%%%%%%%%%%%%%%%%%%%%%%%%%%%%%%%%%%%%%%%%%%%%%%%%%%%%%%%%%
%% AGFA-Book-Einleitung.tex -- Kapitel ohne Nummer
%% \addchap steht im Inhaltsverzeichnis und im Kolumnentitel,
%% bekommt aber keine Nummer.
%% %%%%%%%%%%%%%%%%%%%%%%%%%%%%%%%%%%%%%%%%%%%%%%%%%%%%%%%%%%%%%%%
\addchap{Einleitung}
\label{chap:einleitung}

Diese Vorlage zeigt den Aufbau einer Arbeit mit Kapiteln. In
\cref{chap:grundlagen} werden die Begriffe eingeführt, in
\cref{chap:anwendungen} angewendet. Das Hauptresultat ist
\cref{thm:banach}; es stützt sich auf \cref{lem:cauchy}.

Innerhalb eines Kapitels heißen Sätze und Gleichungen kurz
\enquote{Theorem 2.1} bzw.\ \enquote{(2.1)}: Die erste Zahl ist die
Section, die zweite zählt in der Section. Bei Verweisen wird die
Kapitelnummer vorangestellt, also \enquote{Theorem I.2.1}. So ist jeder
Verweis eindeutig, egal von wo er kommt.
		% Einleitung, ohne Nummer
%% %%%%%%%%%%%%%%%%%%%%%%%%%%%%%%%%%%%%%%%%%%%%%%%%%%%%%%%%%%%%%%%
%% AGFA-Book-Kapitel-1.tex -- Kapitel I
%% %%%%%%%%%%%%%%%%%%%%%%%%%%%%%%%%%%%%%%%%%%%%%%%%%%%%%%%%%%%%%%%
\chapter{Grundlagen}
\label{chap:grundlagen}

Ein kurzer Überblick über das Kapitel steht vor der ersten Section.

\section{Metrische Räume}
\label{sec:metrisch}

\begin{definition}
\label{def:metrik}
Eine Abbildung $d\colon X\times X\to\mathbb{R}$ heißt \emph{Metrik}, wenn
für alle $x,y,z\in X$ gilt
\begin{equation}
\label{eq:dreieck}
  d(x,z)\le d(x,y)+d(y,z).
\end{equation}
\end{definition}

\begin{example}
Der Betrag liefert die Metrik $d(x,y)=\abs{x-y}$ auf $\mathbb{R}$.
\end{example}

\section{Vollständigkeit}
\label{sec:vollstaendig}

\begin{lemma}
\label{lem:cauchy}
Jede konvergente Folge ist eine Cauchy-Folge.
\end{lemma}

\begin{proof}
Folgt aus der Dreiecksungleichung \eqref{eq:dreieck}.
\end{proof}

\subsection{Kontraktionen}
\label{subsec:kontraktion}

\begin{theorem}[Banach]
\label{thm:banach}
Ist $(X,d)$ vollständig und $T\colon X\to X$ eine Kontraktion, so hat $T$
genau einen Fixpunkt.
\end{theorem}

\begin{corollary}
\label{cor:eindeutig}
Die Fixpunktiteration konvergiert für jeden Startwert.
\end{corollary}

Die Abschätzung
\begin{equation}
\label{eq:apriori}
  d(x_n,x^*)\le\frac{q^n}{1-q}\,d(x_1,x_0)
\end{equation}
zeigt die Geschwindigkeit.
		% Kapitel I
%% %%%%%%%%%%%%%%%%%%%%%%%%%%%%%%%%%%%%%%%%%%%%%%%%%%%%%%%%%%%%%%%
%% AGFA-Book-Kapitel-2.tex -- Kapitel II
%% %%%%%%%%%%%%%%%%%%%%%%%%%%%%%%%%%%%%%%%%%%%%%%%%%%%%%%%%%%%%%%%
\chapter{Anwendungen}
\label{chap:anwendungen}

\section{Differentialgleichungen}
\label{sec:dgl}

Wir wenden \cref{thm:banach} aus \cref{subsec:kontraktion} an
(siehe auch \cref{sec:vollstaendig} und \cref{cor:eindeutig}).
Mit \eqref{eq:apriori} erhält man eine Fehlerschranke.

\begin{theorem}[Picard--Lindelöf]
\label{thm:picard}
Ist $f$ Lipschitz-stetig, so hat das Anfangswertproblem genau eine Lösung.
\end{theorem}

\begin{proof}
\Cref{thm:banach} auf den Integraloperator anwenden; die Metrik aus
\cref{def:metrik} ist dabei die Supremumsmetrik.
\end{proof}

\begin{exercise}
\label{ex:picard}
Man führe die Abschätzung in \cref{thm:picard} aus.
\end{exercise}

Verweise auf mehrere Ergebnisse: \cref{lem:cauchy,thm:banach,thm:picard}.
		% Kapitel II
% \include{name}								% etc; bitte den Aufbau beachten

%% -- Anhang: Kapitel heißen dann A, B, ...; Verweise A.1.2
% \appendix
% \include{./content-book/AGFA-Book-Anhang}

%% %%%%%%%%%%%%%%%%%%%%%%%%
%% Backmatter: Kapitel ohne Nummer, Seitenzahlen laufen weiter
%% %%%%%%%%%%%%%%%%%%%%%%%%
\backmatter

%% -- Literaturverzeichnis
%%
%\nocite{}		% für Literatur, die nicht zitiert wurde; 
				% wenig sinnvoll dieses zu tun; siehe ReadMe.pdf
				
%% -- Literatur
%% --	
\RaggedRight
\nocite{voss:2025a,lamport:1986,urysohn:1925a,aigner:2010,axler:2015,lamport:dmv,kahle:2018,zabreiko:1969, comment}
\printbibliography
%% -- Bildquellen (nur wenn \quelle{...} verwendet wurde)
\druckequellen

%% --

\end{document}
